\documentclass{ftpedia}

% --- magazine issue shown in the running header -----------------------------
\ftpediaissue{2026}{3}

% --- article metadata -------------------------------------------------------
\rubrik{Modell}
\titel{Beispielartikel: ein Modell im Maßstab 1:10}
\autor{Max Mustermann}

\begin{document}

\artikelkopf

\begin{ftabstract}
Dies ist ein Beispiel für den typografischen Aufbau eines ft:pedia-Artikels:
Rubrik, Titel, Autor und ein kursive gesetzter Vorspann über die volle
Satzbreite, gefolgt vom zweispaltigen Fließtext.
\end{ftabstract}

\begin{multicols}{2}

\section{Prinzip und Geschichte}

Der Vorläufer des Fotoapparats eignet sich prima als schmuckes Schauobjekt
oder zum perspektivischen Zeichnen. Eine \emph{Camera obscura} ist eine dunkle
Kammer, in die Licht durch ein kleines Loch fällt.\footnote{Eine Fußnote wird
wie gewohnt mit \texttt{\textbackslash footnote} gesetzt.}

Die Brennweite beträgt $f = 103\,\mathrm{mm}$, und für den Winkel gilt
\begin{equation}
  \cos\beta = \frac{a^{2}-b^{2}+c^{2}}{2ac} \tag{1}
\end{equation}

\subsection{1.1 Aufbau}

Als besonders vielseitig erweist sich der Sternmotor. Die Bauteile sind:

\begin{itemize}
  \item Grundbaustein 15,
  \item Drehscheibe 60,
  \item Rastkupplung.
\end{itemize}

\begin{lstlisting}
1 import io
2 while true:
3   if IO.readPin(C_PIN_05):
4     IO.writePin(C_PIN_02, 1)
\end{lstlisting}

In Zeile 3 wird ein Taster abgefragt; der Befehl
\texttt{print(IO.readPin(pin))} gibt den Wert auf dem Terminal aus.

\begin{figure}[H]
  \centering
  \fbox{\parbox[c][35mm][c]{0.85\linewidth}{\centering Abbildungsplatzhalter}}
  \caption{Prinzip der Camera obscura}
\end{figure}

\section{Ergebnisse}

\begin{table}[H]
  \centering
  \begin{tabular}{lrr}
    \hline
    Motor & Zylinder & Leistung \\
    \hline
    SP1  & 1 & 10\,W \\
    V2P  & 2 & 18\,W \\
    \hline
  \end{tabular}
  \caption{Übersicht der Modellvarianten}
\end{table}

Weitere Informationen stehen auf der Projektseite \url{ftpedia.de}.

\end{multicols}

\end{document}
